\documentclass[11pt,a4paper]{article}

% ---------------------------------------------------------------------------
% Relatorio de Perfilamento e Analise
% Projeto: Otimizacao de Telemetria com Buffer Circular
% Compilar com pdflatex (ex.: Overleaf). Usa pgfplots e listings.
% ---------------------------------------------------------------------------
\usepackage[utf8]{inputenc}
\usepackage[T1]{fontenc}
\usepackage[brazil]{babel}
\usepackage{amsmath, amssymb}
\usepackage{geometry}
\geometry{margin=2.5cm}
\usepackage{graphicx}
\usepackage{booktabs}
\usepackage{xcolor}
\usepackage{listings}
\usepackage{pgfplots}
\pgfplotsset{compat=1.17}
\usepackage{hyperref}
\hypersetup{colorlinks=true, linkcolor=blue, urlcolor=blue}

\definecolor{codebg}{rgb}{0.96,0.96,0.96}
\definecolor{codekw}{rgb}{0.1,0.1,0.6}
\definecolor{codecmt}{rgb}{0.3,0.5,0.3}
\lstset{
  basicstyle=\ttfamily\footnotesize,
  backgroundcolor=\color{codebg},
  keywordstyle=\color{codekw}\bfseries,
  commentstyle=\color{codecmt}\itshape,
  language=C++,
  showstringspaces=false,
  breaklines=true,
  frame=single,
  numbers=left,
  numberstyle=\tiny,
}

\title{\textbf{Otimização de Telemetria com Buffer Circular}\\[0.3em]
\large Análise de Algoritmos e Sistemas Embarcados:\\ A Anatomia da Eficiência em Sistemas Embarcados}
\author{Felipe Sérgio \and Thiago Belo \and Thiago Von Sohsten \and Sergio Gouveia \and Enzo Nunes}
\date{2026}

\begin{document}
\maketitle

\begin{abstract}
Este relatório contrasta empiricamente duas estratégias de gerência de uma
janela de amostras de sensores em um sistema embarcado (ESP32) com transmissão
via MQTT: (i) a \textbf{Vertente 1}, anti-padrão baseado em realocação
dinâmica (\texttt{realloc}) e/ou deslocamento de elementos, com custo
$O(n)$ por inserção; e (ii) a \textbf{Vertente 2}, um \emph{Ring Buffer} de
tamanho fixo com índices \emph{head}/\emph{tail} e custo $O(1)$ por operação.
Apresentamos a prova assintótica das complexidades, o diagnóstico do impacto na
saúde do heap e a medição de latência ($\mu s$) sob escalas crescentes de $N$
($100$ a $20\,000$), evidenciando o surgimento de \emph{jitter} na Vertente~1 e
a estabilidade da Vertente~2.
\end{abstract}

\section{Objetivo}
Implementar um sistema de monitoramento para captura de amostras de sensores e
transmissão via protocolo MQTT, contrastando empiricamente a eficiência de uma
abordagem baseada em realocação dinâmica/deslocamento de memória frente a uma
implementação de Buffer Circular de tamanho fixo.

\section{Contextualização do Problema}
Em sistemas embarcados como o ESP32 existe uma disparidade temporal crítica
entre a \emph{captura} de dados (na ordem de microssegundos) e a
\emph{transmissão} de rede (na ordem de milissegundos). Ao usar coleções
dinâmicas ou deslocamento de arrays para manter um histórico de sensores, o
processador consome ciclos lineares ao tamanho da amostra a cada nova leitura.
Isso (a) introduz \emph{jitter} --- instabilidade no período de amostragem ---
e (b) fragmenta o heap a cada \texttt{realloc}, podendo levar à exaustão de
memória e ao colapso do sistema. O Buffer Circular elimina ambos os problemas:
a memória é alocada uma única vez e as operações têm custo constante.

\section{Metodologia: As Duas Vertentes}

\subsection{Vertente 1 --- A Abordagem Ineficiente (Anti-Padrão)}
\begin{itemize}
  \item \textbf{Lógica de Dados:} deslocamento de elementos ($O(n)$) ou uso de
  \texttt{realloc()} a cada nova leitura.
  \item \textbf{Comportamento MQTT:} envio síncrono que bloqueia a amostragem
  durante a latência de rede.
\end{itemize}

\subsection{Vertente 2 --- A Abordagem Eficiente (Buffering Circular)}
\begin{itemize}
  \item \textbf{Lógica de Dados:} \emph{Ring Buffer} com índices \emph{head} e
  \emph{tail}.
  \item \textbf{Complexidade:} inserção e remoção em tempo constante $O(1)$.
  \item \textbf{Comportamento MQTT:} modelo Produtor--Consumidor; o buffer
  absorve a latência da rede sem interromper o sensor.
\end{itemize}

\section{Implementação}
O código-fonte (Entregável~1) está em C++/Arduino no diretório
\texttt{esp32/src/}: a classe do Buffer Circular em \texttt{RingBuffer.h} e o
anti-padrão em \texttt{ShiftBuffer.h} (\texttt{ShiftBuffer} por deslocamento e
\texttt{GrowingBuffer} por \texttt{realloc}). A instrumentação e o benchmark
ficam em \texttt{Telemetry.h}.

\subsection{Ring Buffer ($O(1)$, sem movimentação em massa)}
\begin{lstlisting}
template <typename T, size_t CAPACITY>
class RingBuffer {
public:
  bool push(const T &item) {            // O(1)
    if (isFull()) { tail_ = advance(tail_); count_--; }
    buffer_[head_] = item;
    head_ = advance(head_);
    count_++;
    return true;
  }
  bool pop(T &out) {                    // O(1)
    if (isEmpty()) return false;
    out = buffer_[tail_];
    tail_ = advance(tail_);
    count_--;
    return true;
  }
private:
  size_t advance(size_t i) const { return (i + 1) % CAPACITY; }
  T buffer_[CAPACITY];                  // alocado uma unica vez
  size_t head_, tail_, count_;
};
\end{lstlisting}

\subsection{Anti-padrão (deslocamento $O(n)$ e \texttt{realloc})}
\begin{lstlisting}
// Janela fixa por deslocamento de elementos: O(n) por insercao
void push(const T &item) {
  if (size_ < window_) { data_[size_++] = item; return; }
  for (size_t i = 1; i < window_; i++) data_[i - 1] = data_[i]; // O(n)
  data_[window_ - 1] = item;
}

// Historico crescente via realloc a cada leitura: O(n) por insercao
void push(const T &item) {
  T *grown = (T *)realloc(data_, sizeof(T) * (size_ + 1)); // copia tudo: O(n)
  if (grown == nullptr) return;
  data_ = grown;
  data_[size_++] = item;
}
\end{lstlisting}

\subsection{Instrumentação}
Conforme a seção de instrumentação do enunciado, mede-se tempo de execução e
saúde da memória:
\begin{lstlisting}
unsigned long start = micros();
// Logica de insercao de dados
unsigned long duration = micros() - start;
Serial.printf("Latencia: %lu us | Heap Livre: %u bytes\n",
              duration, ESP.getFreeHeap());
\end{lstlisting}

\section{Análise Assintótica}

\paragraph{Ring Buffer.} As operações \texttt{push} e \texttt{pop} executam um
número fixo de instruções: uma escrita/leitura indexada e uma atualização
modular de índice. Não há laços dependentes de $N$. Logo,
\[
T_{\text{ring}}(\text{push}) = O(1), \qquad
T_{\text{ring}}(\text{pop}) = O(1).
\]
Inserir $N$ amostras custa $\sum_{i=1}^{N} O(1) = O(N)$, com latência
\emph{por inserção} constante (sem jitter).

\paragraph{Deslocamento de janela.} Com a janela cheia, cada \texttt{push}
percorre os $w$ elementos da janela: $O(w)$. Inserir $N$ amostras custa
$O(N \cdot w)$.

\paragraph{Histórico crescente via \texttt{realloc}.} Ao inserir o $i$-ésimo
elemento, o \texttt{realloc}, no pior caso, aloca um novo bloco e copia os $i-1$
elementos existentes: $O(i)$. O custo total de inserir $N$ amostras é
\[
\sum_{i=1}^{N} O(i) = O\!\left(\frac{N(N+1)}{2}\right) = O(N^2).
\]
A latência \emph{por inserção} cresce linearmente com $N$ --- exatamente a
origem do \emph{jitter} observado.

\section{Diagnóstico de Memória}
A Vertente~2 aloca o buffer uma única vez (memória estática de tamanho
\texttt{CAPACITY}); o heap permanece constante e sem fragmentação durante toda a
operação. A Vertente~1, ao chamar \texttt{realloc} a cada leitura, gera um padrão
de alocação/liberação que fragmenta progressivamente o heap: blocos antigos
liberados deixam lacunas que dificilmente são reaproveitadas, reduzindo o maior
bloco contíguo disponível (\emph{watermark}) mesmo quando o total livre ainda é
alto. Em hardware, isso é observável via \texttt{ESP.getFreeHeap()} (queda
progressiva) e via APIs de maior bloco livre; eventualmente o \texttt{realloc}
retorna \texttt{nullptr} e o sistema colapsa por exaustão de memória.

\section{Resultados Experimentais}
A medição abaixo foi obtida com a instrumentação \texttt{micros()} executando o
benchmark de escala (Vertente~1 com \texttt{realloc}/cópia $O(n)$ vs.\
Vertente~2 com Ring Buffer $O(1)$). Cada linha corresponde à inserção de $N$
amostras.

\begin{table}[h]
\centering
\begin{tabular}{rrrr}
\toprule
$N$ & V1 total ($\mu s$) & V2 total ($\mu s$) & Ganho ($\times$) \\
\midrule
100     & 155      & 20  & 7,8 \\
1\,000  & 2\,177   & 80  & 27,2 \\
5\,000  & 23\,128  & 92  & 251,4 \\
20\,000 & 221\,848 & 125 & 1\,774,8 \\
\bottomrule
\end{tabular}
\caption{Latência total por escala. Os números são representativos; reproduza
com \texttt{TELE\_CSV=1 npm run telemetry} (gera \texttt{perf-results.csv}).}
\label{tab:perf}
\end{table}

\begin{figure}[h]
\centering
\begin{tikzpicture}
\begin{axis}[
  width=0.85\textwidth, height=7cm,
  xlabel={$N$ (amostras)}, ylabel={Latência total ($\mu s$)},
  xmode=log, ymode=log, log basis x=10, log basis y=10,
  grid=both, legend pos=north west,
]
\addplot[red, mark=*] coordinates {(100,155)(1000,2177)(5000,23128)(20000,221848)};
\addplot[green!60!black, mark=square*] coordinates {(100,20)(1000,80)(5000,92)(20000,125)};
\legend{Vertente 1 (shift/realloc) $O(N^2)$, Vertente 2 (ring buffer) $O(N)$}
\end{axis}
\end{tikzpicture}
\caption{Escala log--log: a Vertente~1 cresce de forma super-linear ($O(N^2)$)
enquanto a Vertente~2 mantém-se praticamente plana.}
\label{fig:perf}
\end{figure}

A Tabela~\ref{tab:perf} e a Figura~\ref{fig:perf} confirmam o comportamento
previsto na análise de escalabilidade:
\begin{itemize}
  \item $N = 100$: diferença ainda discreta (V1 $\approx 8\times$ mais lenta).
  \item $N = 5\,000$: atraso perceptível na amostragem da V1; a V2 mantém
  latência mínima.
  \item $N = 20\,000$: a V1 leva $\approx 222\,ms$ (alto consumo de CPU e
  \emph{jitter}); a V2 permanece em $\approx 0{,}1\,ms$, sem alteração relevante.
\end{itemize}

\section{Discussão: Instabilidade de Rede}
Sob instabilidade ou gargalo de rede, a publicação MQTT pode levar dezenas a
centenas de milissegundos. Na Vertente~1, com envio \emph{síncrono}, a
amostragem fica bloqueada durante esse intervalo: amostras são perdidas e o
período de captura torna-se irregular. Na Vertente~2, o modelo
\emph{Produtor--Consumidor} desacopla as duas escalas de tempo: o produtor
continua amostrando em $\mu s$ e enfileirando no Ring Buffer em $O(1)$, enquanto
o consumidor drena lotes para a rede quando esta está disponível. O buffer
\emph{absorve} a latência: enquanto não transborda (e, com janela deslizante,
mesmo no transbordo prioriza as amostras mais recentes), a aquisição
permanece estável. Simulamos esse gargalo elevando o intervalo de \emph{flush}
do consumidor e observando que a taxa de amostragem do produtor não se altera.

\section{Conclusão}
A escolha da estrutura de dados é determinante para a viabilidade de um sistema
de telemetria embarcado. A realocação dinâmica/deslocamento ($O(n)$ por
inserção, $O(N^2)$ acumulado) combina jitter crescente e fragmentação de heap,
levando ao colapso sob escala. O Buffer Circular ($O(1)$, memória estática)
entrega latência constante e estabilidade de memória, e seu acoplamento a um
modelo Produtor--Consumidor torna o sistema resiliente à latência de rede.

\section*{Reprodutibilidade}
O painel de telemetria (Entregável~2) está no dashboard React (aba
\emph{Telemetria}). Para reproduzir sem hardware:
\begin{enumerate}
  \item \texttt{cd simulator \&\& npm install}
  \item \texttt{npm run broker} (broker MQTT)
  \item \texttt{npm run telemetry} (benchmark + stream de amostras)
  \item \texttt{cd app/dashboard \&\& npm start} e abra a aba \emph{Telemetria}.
\end{enumerate}
O firmware equivalente (Entregável~1) está em \texttt{esp32/src/} e publica os
mesmos tópicos \texttt{telemetry/*}.

\end{document}
